\documentclass[14 pt]{extarticle}

	\usepackage[french]{babel}
	\usepackage[utf8]{inputenc}  
	\usepackage[T1]{fontenc}
	\usepackage{amssymb}
	\usepackage[mathscr]{euscript}
	\usepackage{stmaryrd}
	\usepackage{amsmath}
	\usepackage{tikz}
	\usepackage[all,cmtip]{xy}
	\usepackage{amsthm}
\usetikzlibrary{arrows.meta}
	\usepackage{varioref}
	\usepackage{geometry}
	\geometry{a4paper}
	\usepackage{lmodern}
	\usepackage{hyperref}
	\usepackage{array}
	 \usepackage{fancyhdr}
	 \usepackage{float}
\renewcommand{\theenumi}{\alph{enumi})}
	\pagestyle{fancy}
	\theoremstyle{plain}
	\fancyfoot[C]{} 
	\fancyhead[L]{Contrôle}
	\fancyhead[R]{16 mars 2026}\geometry{
 a4paper,
 total={190mm,257mm},
 left=10mm,
 top=20mm,
 }
	
	
	\title{Contrôle Chapitre 7}
	\date{}
	\begin{document}

\begin{center}{\Large Contrôle Chapitre 7}\\  
\textbf{On détaillera toujours les calculs.}
 \end{center}
 

\subsection*{Exercice 1 (3 points)}
Calculer :
\begin{enumerate}
\item trois quarts de $28$. 
\item $\frac{7}{12}$ de $144$.
\end{enumerate}
\subsection*{Exercice 2 (2 points)}
\begin{enumerate}
\item Écrire comme un nombre mixte la fraction $\frac{113}{4}$.
\item Écrire comme une seule fraction le nombre mixte $13+ \frac23$.
\end{enumerate}
\subsection*{Exercice 3 (5 points)}
\begin{enumerate}
\item Trouver l'écriture décimale de $\frac{5}8$. 
\item Exprimer la fraction précédente comme un pourcentage.
\item Donner une valeur approchée au centième de $\frac{13}7$. 
\item Donner une valeur approchée au centième de $\frac{4}{39}$. Exprimer également le résultat en pourcentage.
\end{enumerate}

\subsection*{Exercice 4 (4 points)}
Sur la demi-droite graduée suivante : 

\begin{tikzpicture}[scale = .9]
\draw[-{Latex[length=5mm, width=3mm]}](0,0)--(20,0); 
\foreach \x in {1, 2,..., 19}{
\draw (\x, -.2)--(\x, +.2);
}
\foreach \x in {0, 9, 18}{
\draw (\x, -.4)--(\x, +.4);

\draw (\x, -.7) node{\pgfmathparse{\x/9}\pgfmathprintnumber[fixed,precision=0]{\pgfmathresult}};
} 
\foreach \x in {2,  5, 13}{
\draw (\x, 0) -- ++ (45:.3);
\draw (\x, 0) --++ (135:.3);
\draw (\x, 0) --++ (-135:.3);
\draw (\x, 0) --++ (-45:.3);}
\draw (2, .7) node{A};
\draw (5, .7) node{B};
\draw (13, .7) node{C};
\end{tikzpicture}
\begin{enumerate}
\item Placer les points $D\left(\frac59\right)$, $E\left(\frac{17}9\right)$, $F\left(\frac13\right)$ et $G\left(\frac{16}{24}\right)$.
\item Donner les abscisses des points $A$, $B$ et $C$.
\end{enumerate}  

\subsection*{Exercice 5 (6 points)}
 

\begin{enumerate}
\item Écrire une fraction égale à $\frac34$ ayant $48$ comme dénominateur.
\item Écrire une fraction égale à $\frac46$ ayant $21$ comme dénominateur.
\item Les fractions $\frac{12}{31}$ et $\frac{9}{20}$ sont-elles égales ?
\item Simplifier le plus possible la fraction $\frac{210}{315}$.
\item Les fractions $\frac{188}{292}$ et $\frac{24}{36}$ sont-elles égales ?
\end{enumerate}
\newpage

\begin{center}{\Large Contrôle Chapitre 7}\\ 
\textbf{On détaillera toujours les calculs.}
 \end{center}
 

\subsection*{Exercice 1 (3 points)}
Calculer :
\begin{enumerate}
\item trois quarts de $24$. 
\item $\frac{5}{12}$ de $132$.
\end{enumerate}
\subsection*{Exercice 2 (2 points)}
\begin{enumerate}
\item Écrire comme un nombre mixte la fraction $\frac{83}{4}$.
\item Écrire comme une seule fraction le nombre mixte $11+ \frac13$.
\end{enumerate}
\subsection*{Exercice 3 (5 points)}
\begin{enumerate}
\item Trouver l'écriture décimale de $\frac{3}8$. 
\item Exprimer la fraction précédente comme un pourcentage.
\item Donner une valeur approchée au centième de $\frac{11}7$. 
\item Donner une valeur approchée au centième de $\frac{2}{39}$. Exprimer également le résultat en pourcentage.
\end{enumerate}

\subsection*{Exercice 4 (4 points)}
Sur la demi-droite graduée suivante : 

\begin{tikzpicture}[scale = .9]
\draw[-{Latex[length=5mm, width=3mm]}](0,0)--(20,0); 
\foreach \x in {1, 2,..., 19}{
\draw (\x, -.2)--(\x, +.2);
}
\foreach \x in {0, 9, 18}{
\draw (\x, -.4)--(\x, +.4);

\draw (\x, -.7) node{\pgfmathparse{\x/9}\pgfmathprintnumber[fixed,precision=0]{\pgfmathresult}};
} 
\foreach \x in {3, 6, 13}{
\draw (\x, 0) -- ++ (45:.3);
\draw (\x, 0) --++ (135:.3);
\draw (\x, 0) --++ (-135:.3);
\draw (\x, 0) --++ (-45:.3);}
\draw (3, .7) node{A};
\draw (6, .7) node{B};
\draw (13, .7) node{C};
\end{tikzpicture}
\begin{enumerate}
\item Placer les points $D\left(\frac49\right)$, $E\left(\frac{16}9\right)$, $F\left(\frac23\right)$ et $G\left(\frac{10}{15}\right)$.
\item Donner les abscisses des points $A$, $B$ et $C$.
\end{enumerate}  

\subsection*{Exercice 5 (6 points)}
 

\begin{enumerate}
\item Écrire une fraction égale à $\frac34$ ayant $48$ comme dénominateur.
\item Écrire une fraction égale à $\frac46$ ayant $21$ comme dénominateur.
\item Les fractions $\frac{12}{31}$ et $\frac{9}{20}$ sont-elles égales ?
\item Simplifier le plus possible la fraction $\frac{210}{525}$.
\item Les fractions $\frac{182}{273}$ et $\frac{26}{39}$ sont-elles égales ?
\end{enumerate}

 	\end{document}
